This section reviews existing research on ClimateNet, focusing on prior attempts to implement automated detection for atmospheric features. We further examine current literature on Parameter-Efficient Fine-Tuning (PEFT) techniques, which are used to adapt large-scale foundation models for specific downstream tasks. Additionally, we discuss recent advancements in SAM adaptation, particularly focusing on methods that aim to transition SAM toward a prompt-free architecture for automated segmentation
\subsection{Automated detection of atmospheric features}
Traditional atmospheric feature detection relies on rule-based algorithms that use physical and mathematical heuristics, such as identifying Tropical Cyclones through minima or maxima in sea level pressure \cite{bourdin2022intercomparison, neu2013imilast}. Atmospheric fronts are typically identified using derivatives of thermal variables and threshold-based filters, while Atmospheric Rivers are detected based on specific moisture thresholding and geometric requirements\cite{beckert2023three, neu2013imilast}. However, these rule-based systems are often inconsistent across different studies and can be difficult to formulate into a universal set of physical parameters.

Recent research demonstrates that supervised learning with convolutional neural networks, such as U-Net and CG-Net, can learn these detection tasks as semantic segmentation problems \cite{lagerquist2019deep, biard2019automated}. Using the ClimateNet dataset, these models achieve performance comparable to human experts by classifying grid points into background, TC, or AR classes \cite{kapp2020spatio, prabhat2020climatenet}. Using CNNs for feature detection via semantic segmentation offers several key advantages, including significantly increased computational performance \cite{neu2013imilast, higgins2023using, radke2024explaining}. Furthermore, CNNs provide the unique capability to learn complex atmospheric patterns that are otherwise difficult to explicitly formulate as a rigid set of physical or mathematical rules \cite{niebler2022automated, tian2023deep, radke2024explaining}.

Findings using Layer-wise Relevance Propagation (LRP) confirm that these CNNs learn physically plausible patterns; for TCs, the models prioritize point-shaped extrema in total precipitable water and circular wind motion, while for ARs, they focus on elongated bands of high humidity and eastward winds. Additionally, studies show that while z-score normalization can cause an ``ignorant-to-zero-input" issue in explainability maps, shifting the data ensures the models correctly utilize information across the entire input range to inform their spatial decisions \cite{radke2024explaining}.
\subsection{Parameter-Efficient Tuning}
Parameter-efficient tuning for foundation models has gained significant importance as a means to address the challenges of massive model scales and the inefficiency of maintaining unique versions for every task. Typically, the backbone parameters, such as the Vision Transformer in SAM, remain frozen, while a small set of auxiliary trainable weights is introduced and adapted to the specific task.  

One way to implement these lightweight parameters is prompt tuning, which involves adding learnable tokens to the input sequence at various transformer layers to act as instructions that steer the frozen model. For instance, Visual Prompt Tuning  introduces a small number of trainable parameters into the input space\cite{jia2022visual}, demonstrating that tuning less than 1\% of the model can outperform full fine-tuning. Similarly, MaPLe expands this by learning separate prompts across different stages of both vision and language branches to progressively model stage-wise feature relationships\cite{khattak2023maple}, while Context Optimization models context words as learnable continuous vectors to adapt models like CLIP for downstream recognition \cite{zhou2022learning}.

Another effective approach to perform for lightweight fine-tuning is by the integration of learnable sub-networks directly into the transformer modules. These adapters typically consist of small bottleneck layers that learn task-specific refinements without modifying the original weights. For example, CLIP-Adapter employs two-layer MLP modules to fine-tune high-level features of frozen vision-language models \cite{gao2025clipadapterbettervisionlanguagemodels}. Residual adapters offer a parallel approach by adding tunable residual modules that steer a frozen network toward diverse visual domains\cite{rebuffi2017learningmultiplevisualdomains}. More specialized modules include Explicit Visual Prompting, which introduces a tuner focusing on high-frequency components to segment low-level structures\cite{liu2023explicitvisualpromptinglowlevel}, and the Side Adapter Network (SAN), which attaches a lightweight side network to predict mask proposals for open-vocabulary tasks\cite{xu2023adapternetworkopenvocabularysemantic}. While these methods introduce additional architectural components alongside the frozen backbone, a distinct category of parameter-efficient tuning focuses on the mathematical re-parameterization of the weight updates themselves. This approach is exemplified by low-rank adaptation techniques, which modify the internal dynamics of the model without increasing its depth or inference latency.

\subsubsection{LoRA}
The development of Low-Rank Adaptation (LoRA) addresses the computational challenges of fine-tuning large-scale foundation models by assuming that task-specific weight updates possess a low intrinsic rank. Instead of updating the entire pre-trained weight matrix $W_0 \in \mathbb{R}^{d \times k}$, LoRA represents the update $\Delta W$ as a product of two lower-rank matrices $B \in \mathbb{R}^{d \times r}$ and $A \in \mathbb{R}^{r \times k}$, where $r \ll \min(d, k)$. During training, the original backbone weights remain frozen, and only these low-rank decomposition matrices are updated, which drastically reduces the number of trainable parameters. This strategy has been proven effective in various domains, such as the MeLo framework for medical diagnosis, which achieved competitive results using only 0.17\% of the original parameters \cite{zhu2024melo}. Recent adaptations of SAM have successfully utilized LoRA to bridge the domain gap between natural and specialized datasets, such as medical or atmospheric imagery, while preventing catastrophic forgetting \cite{wu2025multi, zhang2024improving}.

\subsubsection{Conditional joint tuning strategy}
Previous studies often focused on tuning only a single backbone or dual-encoder systems (like CLIP). Because SAM has a disproportionately small mask decoder compared to its heavyweight image encoder, simply adding prompts or adapters to one part might not effectively bridge the gap for specialized data. To address these structural challenges, recent frameworks like HQ-SAM \cite{ke2023segment} and CAT-SAM \cite{xiao2024cat} employ a conditional joint tuning strategy that prioritizes the relationship between the encoder and decoder. HQ-SAM (High-Quality SAM) first addresses the limitation of coarse mask generation by introducing a small learnable "HQ-Output" token into the mask decoder. This token incorporates global image context and latent mask information, which is then combined with the original embeddings to facilitate high-quality mask generation for complex, domain-specific structures. Building upon this concept of decoder-based refinement, CAT-SAM introduces a prompt bridge structure to facilitate decoder-conditioned joint tuning. This mechanism explicitly maps domain-specific features from the mask decoder back into the heavyweight image encoder, fostering a synergy that allows both components to adapt together in a data-efficient manner.

\subsection{Prompt Generator for SAM}


To help SAM work without someone manually clicking or drawing boxes, research has shifted toward making the model prompt-free. The manual prompting is a major bottleneck for real-time tasks like weather forecasting, industrial automation, or medical imaging, where you might need to segment thousands of objects instantly. While vanilla SAM attempts ``unintelligent automation" through dense grid searches \cite{kirillov2023segment}, this method often misses small objects or generates redundant masks, slowing down processing. Alternative approaches utilize another detection models like YOLOv8 to generate auxiliary masks, but these external, heavyweight architectures introduce significant computational overhead and lack tight alignment with SAM’s internal features \cite{zhang2023personalize}.

More sophisticated strategies focus on self-prompting and iterative refinement, where SAM generates its own guidance. For instance, SAM-SP uses initial vanilla predictions to create bounding boxes for subsequent refinement \cite{zhou2024sam}, while Self-Prompt SAM integrates a multi-scale prompt generator (MSPGenerator) with the image encoder to predict auxiliary masks\cite{xie2025self}. These masks are further processed via Euclidean distance transforms to extract optimal, centralized points and boxes as prompts. To integrate automation natively, lightweight prediction networks like those in AoP-SAM operate directly on image embeddings to produce confidence maps, utilizing adaptive sampling and filtering to eliminate redundant prompts \cite{chen2025aop}. Similarly, UN-SAM fuses multi-scale embeddings to predict high-confidence foreground regions, while the SAN (Segment Any Nuclei) framework learns spatial priors of cell nuclei to automatically generate the thousands of prompts needed for medical imaging\cite{chen2024sam}.

Finally, for highly specialized domains, automated prompting is enhanced by integrating expert knowledge or domain-adaptive mechanisms. PG-SAM leverages clinical diagnostic reports encoded through Vision-Language Models to create adaptive text embeddings that pinpoint high-morbidity regions for precise spatial prompting \cite{wu2025multi}. Meanwhile, DAPSAM addresses domain shifts by using a prototype-based generator powered by a parameterized memory bank \cite{wei2024prompting}. By matching instance-level prototypes from target embeddings against stored source-domain knowledge, the system yields robust, domain-adaptive prompts that autonomously guide SAM’s mask decoder.




